\documentclass[aspectratio=169]{beamer}

% Shared Beamer settings for Elements of AIML lectures (UPES).
% Keep this file free of \documentclass to allow reuse via \input.

% Allow lecture sources to use shared theme/assets without copying them into each lecture folder.
% NOTE: These paths are relative to `Unit-xx/Lecture-yy/latex/` (and `Lab/Experiment-xx/latex/`).
\makeatletter
\def\input@path{{../../../_shared/}{../../../_shared/UPES-Beamer-Theme/}}
\makeatother

\usetheme{UPES}
\setbeamertemplate{navigation symbols}{}

\usepackage{amsmath}
\usepackage{amssymb}
\usepackage{booktabs}
\usepackage{graphicx}
\usepackage{hyperref}
\usepackage{xcolor}
\usepackage{listings}

% Code listings (general-purpose).
\lstset{
  basicstyle=\ttfamily\small,
  keywordstyle=\color{blue},
  commentstyle=\color{gray},
  stringstyle=\color{teal},
  showstringspaces=false,
  columns=fullflexible,
  breaklines=true
}

\usepackage{tikz}
\usetikzlibrary{arrows.meta,positioning,shapes.geometric,calc}

% Per-slide source line (references shown on the slide itself).
\newcommand{\slidesources}[1]{%
  \vfill
  {\tiny\color{gray}\raggedright\sloppy\parbox{\linewidth}{Sources: #1}\par}%
}

% Unit-wise source strings for this subject (used by slides/notes).
% Path is relative to the lecture `latex/` folder (the compilation CWD).
% Source strings for Elements of AIML (used by slides + notes).

\newcommand{\AIMLRepoURL}{https://github.com/tali7c/Elements-of-AIML}

\newcommand{\AIMLLabSources}{%
Provost and Fawcett, \textit{Data Science for Business}.;%
McKinney, \textit{Python for Data Analysis}.;%
WEKA Documentation (Waikato Environment for Knowledge Analysis), accessed 2026-02-28.;%
Matplotlib and Seaborn documentation, accessed 2026-02-28.%
}

\newcommand{\AIMLUnitISources}{%
Rich and Knight, \textit{Artificial Intelligence}, McGraw Hill, 2017.;%
Russell and Norvig, \textit{Artificial Intelligence: A Modern Approach} (4e), Ch. 1--2 (Introduction, Intelligent Agents).;%
Poole and Mackworth, \textit{Artificial Intelligence: Foundations of Computational Agents} (2e), selected sections.%
}

\newcommand{\AIMLUnitIISources}{%
Russell and Norvig, \textit{Artificial Intelligence: A Modern Approach} (4e), Ch. 7--9 (Logical Agents, First-Order Logic, Inference).;%
Huth and Ryan, \textit{Logic in Computer Science} (2e), selected sections on propositional and first-order logic.;%
Genesereth and Nilsson, \textit{Logical Foundations of Artificial Intelligence}, selected chapters.%
}

\newcommand{\AIMLUnitIIISources}{%
Mueller and Massaron, \textit{Machine Learning for Dummies}, For Dummies, 2016.;%
Mitchell, \textit{Machine Learning}, selected chapters on learning basics and evaluation.;%
Scikit-learn documentation, accessed 2026-02-28.%
}

\newcommand{\AIMLUnitIVSources}{%
Mueller and Massaron, \textit{Machine Learning for Dummies}, For Dummies, 2016.;%
James, Witten, Hastie, and Tibshirani, \textit{An Introduction to Statistical Learning} (2e), selected chapters.;%
Scikit-learn documentation, accessed 2026-02-28.%
}

\newcommand{\AIMLUnitVSources}{%
NIST, \textit{AI Risk Management Framework (AI RMF 1.0)}, 2023.;%
Barocas, Hardt, and Narayanan, \textit{Fairness and Machine Learning} (online book), selected sections.;%
Knaflic, \textit{Storytelling with Data}, selected chapters on communicating results.%
}



\title[Elements of AIML]{Elements of AIML}
\subtitle{Unit 04 -- Lecture 02: Supervised Learning (Classification)}
\author{Tofik Ali}
\institute{School of Computer Science, UPES Dehradun}
\date{\today}

\graphicspath{{../images/}{../../../_shared/}{../../../_shared/UPES-Beamer-Theme/}}

\begin{document}

\UPESCoverFrame
\setcounter{framenumber}{0}

\begin{frame}
  \titlepage
  \vspace{0.3em}
  \begin{center}
    \footnotesize Topic focus: predicting \textbf{class labels} and evaluating with proper metrics\\[0.25em]
    \footnotesize Repository: \texttt{\AIMLRepoURL}
  \end{center}
  \slidesources{\AIMLUnitIVSources}
\end{frame}

\begin{frame}{Quick Links}
  \centering
  \hyperlink{sec:cls}{\beamerbutton{Classification}}\hspace{0.6em}
  \hyperlink{sec:cm}{\beamerbutton{Confusion}}\hspace{0.6em}
  \hyperlink{sec:metrics}{\beamerbutton{Metrics}}\hspace{0.6em}
  \hyperlink{sec:imb}{\beamerbutton{Imbalance}}\hspace{0.6em}
  \hyperlink{sec:summary}{\beamerbutton{Summary}}
\end{frame}

\begin{frame}{Agenda}
  \tableofcontents
\end{frame}

\section{Classification basics}
\label{sec:cls}

\begin{frame}{Learning Outcomes}
  By the end of this lecture, you should be able to:
  \begin{itemize}[<+->]
    \item Define classification and give examples.
    \item Understand the confusion matrix and the meaning of TP/FP/TN/FN.
    \item Compute and interpret accuracy, precision, recall, and F1.
    \item Explain thresholds, ROC-AUC intuition, and why imbalance matters.
  \end{itemize}
\end{frame}

\begin{frame}{Abbreviations \& Symbols}
  \begin{itemize}[<+->]
    \item TP/FP/TN/FN: true/false positives/negatives
    \item Precision $=TP/(TP+FP)$,\; Recall/TPR $=TP/(TP+FN)$
    \item FPR $=FP/(FP+TN)$,\; F1: harmonic mean of precision and recall
    \item ROC: TPR vs FPR,\; AUC: area under curve
    \item Threshold $\tau$: score $\ge \tau \Rightarrow$ predict positive
  \end{itemize}
\end{frame}

\begin{frame}{What is Classification?}
  \begin{block}{Classification}
    Predict a \textbf{discrete class label} (binary or multi-class).
  \end{block}
  Examples:
  \begin{itemize}[<+->]
    \item spam vs not spam
    \item disease vs healthy
    \item handwritten digit (0--9)
  \end{itemize}
\end{frame}

\begin{frame}{Common Classification Models (Overview)}
  \begin{itemize}[<+->]
    \item \textbf{Logistic regression}: linear decision boundary + probabilities
    \item \textbf{kNN}: classify by nearest neighbors (distance-based)
    \item \textbf{Decision trees / Random forests}: rule-like splits, non-linear
    \item \textbf{Naive Bayes}: probabilistic model (works well for text)
    \item \textbf{SVM}: margin-based classifier (good with clear separation)
  \end{itemize}
  \vspace{0.3em}
  \small Model choice depends on data size, feature type, interpretability, and accuracy needs.
\end{frame}

\section{Confusion matrix}
\label{sec:cm}

\begin{frame}{Confusion Matrix (Binary)}
  \centering
  \begin{tabular}{c|cc}
    & Pred + & Pred - \\\hline
    True + & TP & FN \\
    True - & FP & TN \\
  \end{tabular}
  \vspace{0.4em}
  \begin{itemize}
    \item TP: correct positive, FP: false alarm
    \item FN: missed positive, TN: correct negative
  \end{itemize}
\end{frame}

\begin{frame}{Worked Example (Compute Metrics)}
  Suppose: TP=40,\; FP=10,\; FN=20,\; TN=930 (rare positives).
  \begin{itemize}[<+->]
    \item Accuracy $= (TP+TN)/(TP+TN+FP+FN)$
    \item Precision $= TP/(TP+FP)=40/50=0.80$
    \item Recall $= TP/(TP+FN)=40/60\approx 0.67$
    \item F1 $= 2\cdot \frac{0.80\cdot 0.67}{0.80+0.67}\approx 0.73$
  \end{itemize}
  \vspace{0.3em}
  \small High accuracy is easy with many TNs; precision/recall reveal minority-class performance.
\end{frame}

\section{Metrics}
\label{sec:metrics}

\begin{frame}{Key Metrics}
  \begin{itemize}[<+->]
    \item Accuracy: $(TP+TN)/(TP+TN+FP+FN)$
    \item Precision: $TP/(TP+FP)$ (how many predicted positives are correct?)
    \item Recall: $TP/(TP+FN)$ (how many true positives did we find?)
    \item F1: harmonic mean of precision and recall
  \end{itemize}
\end{frame}

\begin{frame}{Thresholds and ROC-AUC (Intuition)}
  \begin{itemize}[<+->]
    \item Many classifiers output a score/probability.
    \item A \textbf{threshold} converts a score into a class.
    \item ROC curve shows TPR vs FPR as threshold changes.
    \item AUC summarizes ranking quality (higher is better).
  \end{itemize}
\end{frame}

\begin{frame}{ROC vs PR (When imbalance is high)}
  \begin{itemize}[<+->]
    \item ROC curve can look good even when positives are rare.
    \item Precision-Recall (PR) curve focuses on the positive class.
    \item For imbalanced problems, PR-AUC is often more informative than ROC-AUC.
  \end{itemize}
\end{frame}

\section{Imbalance}
\label{sec:imb}

\begin{frame}{Imbalanced Data: Why Accuracy Fails}
  \begin{exampleblock}{Example}
    If 99\% emails are not spam, a classifier that always predicts ``not spam'' gets 99\% accuracy.
  \end{exampleblock}
  \vspace{0.4em}
  \begin{itemize}[<+->]
    \item Accuracy can be misleading in imbalance.
    \item Use precision/recall, F1, and domain-specific costs.
    \item Lab Experiment-09 focuses on imbalance handling techniques.
  \end{itemize}
\end{frame}

\section{Summary}
\label{sec:summary}

\begin{frame}{Summary}
  \begin{itemize}
    \item Classification predicts labels; evaluate with confusion matrix metrics.
    \item Precision/recall trade-offs depend on error costs.
    \item Threshold choice matters; ROC-AUC provides a ranking view.
    \item Imbalance requires careful metrics and techniques.
  \end{itemize}
  \slidesources{\AIMLUnitIVSources}
\end{frame}

\begin{frame}{Exit Questions}
  \begin{enumerate}
    \item Define TP, FP, TN, FN with a medical test example.
    \item Why is accuracy misleading on imbalanced datasets?
    \item When is recall more important than precision (give one scenario)?
    \item What is the role of a classification threshold?
  \end{enumerate}
\end{frame}

\end{document}
